
\noindent \uline{Lorsqu'on considère les deux membres de} (7.3.1) \uline{comme des modules sur} @@M0@@ \uline{grâce aux homomorphismes d'anneaux} @@M1@@ \uline{et} @@M2@@ (IV 2.7 b) \uline{et aux structures de modules de} @@M3@@ \uline{et} @@M4@@ \uline{sur} @@M5@@ \uline{et} @@M6@@ \uline{respectivement} (IV 2.10), \uline{l'homomorphisme} @@M7@@ \uline{est un homomorphisme de} @@M8@@-\uline{modules}.

La dernière assertion provient du fait que @@M9@@ est un homomorphisme de @@M10@@-modules, ainsi que l'épimorphisme @@M11@@, en vertu de (IV 2.12.2).

Nous allons passer en revue d'autres propriétés du morphisme de spécialisation.

\subsection*{7.4. Compatibilité aux filtrations.}
Munissons les @@M12@@ de la filtration @@M13@@ par la dimension des supports (1.1.1). Alors, si @@M14@@ est de type fini, plus généralement si @@M15@@ est caténaire, \emph{l'homomorphisme de spécialisation est compatible avec ces filtrations}, et induit par suite un homomorphisme (appelé également homomorphisme de spécialisation) :
@@M16@@

Pour le voir, on est ramené par 1.1.4 à prouver que si @@M17@@ est un sous-schéma intègre de dimension @@M18@@ de @@M19@@, alors @@M20@@ est de filtration @@M21@@. Or soit @@M22@@ l'adhérence schématique de @@M23@@ dans @@M24@@, de sorte que @@M25@@ est plat sur @@M26@@ (EGA IV 2.8.1). Il en résulte aisément que @@M27@@, où les deux @@M28@@ sont pris respectivement dans @@M29@@ et dans @@M30@@, d'où @@M31@@. On vérifie aisément (grâce au fait que @@M32@@ est caténaire, donc @@M33@@ caténaire) que @@M34@@ ; dans le cas où @@M35@@ est de type fini, cela résulte aussitôt, par réduction au cas @@M36@@ intègre, de EGA IV 7.1.13. D'où la conclusion annoncée.

